% ============================================================
%  experimental_design.tex  —  §5
%  Updated: 29/Sep/2026
% ============================================================

\section{Experimental design}
\label{sec:design}

\subsection{Three entry points for the weight vector}

A dominance test involves weights at three distinct places, and the
distinction is not cosmetic:
\begin{enumerate}[label=(\alph*)]
  \item $\wsc$ --- the weights in the \emph{method score}, which order
    the sites and hence determine the assembled portfolio $t=M(\wsc)$;
  \item $\wob$ --- the weights used to map the assembled portfolio to
    objective space, $\mathbf{F}(t;\wob)$;
  \item $\wfr$ --- the weights used to build the reference set,
    $\Fset(\wfr)$ of~\eqref{eq:front}.
\end{enumerate}

\begin{remark}[coherence condition]
\label{rem:coherence}
A dominance test is well defined only if $\wob=\wfr$.
If they differ, the portfolio and the members of the reference set are
points of different objective spaces and their componentwise comparison
has no decision-theoretic content.
We state this explicitly because it is easy to violate silently when the
reference set is precomputed once and the Monte Carlo loop recomputes
portfolio objectives per draw.
\end{remark}

Weights are perturbed as $w\sim\mathrm{Dir}(\kappa\,\wahp)$ with
concentration $\kappa\in\{50,200\}$ and $\ndraws$ draws per cell; larger
$\kappa$ concentrates the draws around $\wahp$.
Throughout, $\kappa$ denotes only this concentration parameter; the
separation threshold is always written $d_{\min}=\dmin$~km.

\subsection{Two robustness measures}

Both quantities below satisfy the coherence condition of
Remark~\ref{rem:coherence}:
\begin{align}
  Q_1 &= \Pr_{w}\bigl[\,M(w)\in\Fset(\wahp)\,\bigr],
    && \wsc=w,\quad \wob=\wfr=\wahp, \label{eq:Q1} \\
  Q_2 &= \Pr_{w}\bigl[\,M(w)\in\Fset(w)\,\bigr],
    && \wsc=\wob=\wfr=w. \label{eq:Q2}
\end{align}
They answer different questions.
$Q_2$ is \emph{internal consistency}: given a weight vector, does the
method recover a portfolio that is Pareto-efficient under that same weight
vector?
$Q_1$ is \emph{robustness to weight misspecification}: how often does a
perturbed method still land on the reference set of the nominal decision
maker?
$Q_2$ is the primary measure reported in this work.
$Q_1$ is reported as absolute levels with confidence intervals and is not
differenced across methods; because $Q_1$ with min-max normalisation
evaluates portfolios chosen under various normalisation methods using a
reference set built with min-max, any contrast between methods in $Q_1$
conflates normalisation effects with evaluator alignment.

\begin{proposition}[additive consistency]
\label{prop:additive}
For any reference set choice $\wfr=w$, the optimal-assembly SAW portfolio
under $\wsc=w$ with the same normalisation is non-dominated in $\Fset(w)$.
Consequently $Q_2=1$ for that cell.
\end{proposition}
\noindent
This is the standard weighted-sum scalarisation argument: maximising
$\sum_\ell \lambda_\ell F_\ell$ with $\lambda>0$ yields a Pareto-optimal
point.
We use it as a \emph{regression test} on the pipeline rather than as a
contribution: any implementation returning $Q_2<1$ for SAW optimal under
its own normalisation is faulty.
It also delimits the scope of the factorial contrast, since one of the
eight cells is fixed at unity by construction.

\subsection{Normalisation reference set}

Min-max and vector constants may be computed over all $\Ncand$ candidates
(space~$A$) or over the $\Nvia$ admissible ones (space~$A'$).
Space~$A$ reproduces the source study ($|\Fset|=\FAsize$ portfolios);
space~$A'$ is the factorial design ($|\Fset|=\Fmmsize$), and the two
reference sets share only $\FAinter$ portfolios.
Including inadmissible sites shifts the scale without those sites being
selectable --- a set-dependence mechanism documented for TOPSIS by
\citet{GarciaCascales2012}.

\subsection{Factorial design}

The main experiment is a balanced $2^3$: aggregation
$\{\text{SAW},\text{TOPSIS}\}$ $\times$ normalisation
$\{\text{min--max},\text{vector}\}$ $\times$ assembly
$\{\text{greedy},\text{optimal}\}$, evaluated in $Q_2[A']$.
A wider 12-cell design adds VIKOR, PROMETHEE~II, and rank normalisation;
it is unbalanced (VIKOR under vector normalisation is redundant by affine
invariance) and is reported only for marginal means.
Effects are reported in percentage points with $\nboot$-replicate
bootstrap intervals, paired by draw.

\subsection{Evaluator robustness}
\label{sec:evaluators}

Because the reference set of the $2^3$ design uses min-max partial value
functions, part of the normalisation effect may reflect alignment with the
evaluator rather than an intrinsic property.
We therefore rebuild the reference set under two alternative partial value
functions, holding $\Tset$ fixed:
\begin{itemize}
  \item $\Fset_{\mathrm{vec}}$: vector-normalised partial values
    ($|\Fset|=\Fvecsize$ portfolios);
  \item $\Fset_{\mathrm{fix}}$: partial values on \emph{a priori} fixed
    bounds, independent of the candidate sample
    ($|\Fset|=\Ffixsize$ portfolios).
\end{itemize}
The second follows the recommendation of \citet{GarciaCascales2012} that
the value scale should not depend on the alternative set.
The bounds are conventional physical and planning ranges fixed before the
analysis; we make no regulatory claim for them, since the applicable
jurisdiction is Cape Verde and the fixed-bound evaluator serves only to
break sample dependence, which requires independence from the sample, not
legal authority.
\todo{D2/slope: if any bound is to be presented as normative, convert and
  cite correctly --- a 30\% gradient is $\arctan(0.30)\approx16.7^\circ$,
  not $20^\circ$.}

In this dataset, the two alternative evaluators produce nearly identical
classifications: the reference sets $\Fset_{\mathrm{vec}}$ and
$\Fset_{\mathrm{fix}}$ share only 29 portfolios in common
($|\Fset_{\mathrm{vec}}\cap\Fset_{\mathrm{fix}}|=29$ out of
$\Fvecsize$ and $\Ffixsize$ respectively), yet the
portfolio-level dominance verdict is the same under both evaluators in
99.9\% of the \ndraws{} Monte Carlo draws (Section~\ref{sec:results}).
This coincidence indicates that, for the Fogo Island instance, the
distinction between a vector-normalised and a sample-independent evaluator
is not empirically consequential: both yield essentially the same
partition of the portfolio space into dominated and non-dominated sets.
The result is reported for completeness and as a dataset-specific
robustness check, not as a claim about external evaluators in general.
